\section{De un soporte fijo al alcance logarítmico}
\paragraph{Enunciado relativo a las entradas del acta de jets.}
Fíjense los márgenes de difusión y brecha. Para un soporte físico pertinente $U$ de $v$ sitios,
sea $d_i$ el máximo tamaño en $U$ de una fibra que contiene $i$, y $c$ el número de componentes
de su grafo de incidencia completo, con filas y columnas. Se usa la proyección alternada
sobre las caras de $\log\widehat Z-\Lambda$, con el mismo exterior, campo afín y cofactor.
La denotamos por $H_q(s,\zeta)$; $s=1/j$ sólo en el valor físico.
Bajo las identidades locales P03--05, U06--07 y Z01--02, y los mayorantes de G02--04 y B02,
las anulaciones y divisibilidades locales valen en el polidisco de clase para $q=4v+12$.
Para $j=\lfloor bn\rfloor$, $b>0$ fijo y $5\le v\le8\log n$, se obtiene
\[
 |H_q(1/j,\zeta^{\rm fis})|\le C^v(Cv)^{c+1}j^{-p}\prod_{i\in U}d_i .
\]
El exponente $p$ es $v+c+1$, $v+c$ o $v+c-1$ según haya cero, uno o varios ciclos.
En esta parte se demuestra el paso entre los dos alcances. Las entradas que determinan
el germen y su radio son G02:55--56, B02:43--50 y U02:54--55:
la unicidad de la silla y la misma rama logarítmica identifican el germen.

\paragraph{Truncación sobre una misma rama.}
Se fija primero $U$, con sus datos balanceados. En X03, líneas 2--24, $V_l$ es divisible
por $u^l$, y los momentos centrados satisfacen
\[
 (1-iu)\mu_{i+1}=iu(2\mu_i+\mu_{i-1}),\quad \mu_0=1,\quad\mu_1=0 .
\]
Por inducción su orden es al menos $\lceil i/2\rceil$. Para $Q\ge q$, truncar el exceso
en $q$ y los momentos en $2q+1$ cambia la suma y su logaritmo en el ideal $(u^{q+1})$:
las derivadas en la variable de la mezcla conservan el factor y $d(1)\ne0$ fija el logaritmo.
Al sustituir $u=s/(1-\theta_0)$, el denominador es una unidad. Se sustituyen simultáneamente
los parámetros holomorfos en las caras, se suman átomos y se toma el logaritmo con el mismo
valor basal invertible. Con la truncación de Stirling y una división por $s$ se conserva
al menos el ideal $(s^{q-1})$ (X03, líneas 26--34). La congruencia es de gérmenes
conjuntos en $s,\zeta$, y por ello admite restricciones a caras y derivadas en los campos.
Para leer la identidad algebraica se toma un orden $Q$ suficientemente alto, con $Q\ge q$.
No se compara con un corte antiguo menor que el grado leído. Todos los grados usados
son a lo sumo $c+1\le v$ en estos soportes: cada componente pertinente tiene al menos dos sitios.
Así $q-1>v$ basta. Las constantes de una estimación del aproximante de orden $Q$
no intervienen en este cambio de coeficientes.

\paragraph{Continuación en las activaciones.}
El objeto de orden físico $q$ es holomorfo en $|s|<\kappa/v$,
$|\zeta_i|<\kappa'/d_i$. Cada logaritmo centrado por cara vale $O(v)$ allí
y la suma alternada está acotada por $C^v$, reduciendo ambos radios por un factor fijo.
Esta entrada procede de G02--04 y B02. Fijado $l$, la integral de Cauchy en $s$
muestra que $f_l(\zeta)=[s^l]H_q$ es holomorfa en todo el polidisco de campos.
Si el germen $f_l$ es cero en un entorno abierto de las activaciones independientes,
el teorema de identidad da $f_l=0$ en el polidisco conexo.
El radio del entorno donde se probó la identidad puede depender de $U$.

La proyección por caras anula $H_q$ en cada $\zeta_i=0$, pues los valores se identifican
allí con el mismo cofactor. En la clase unicíclica fíjese su ciclo $\mathcal C$.
En grado $c-1$ el germen proyectado pertenece al ideal
\[
 \left(\prod_{i\in U}\zeta_i\prod_{i\in\mathcal C}\zeta_i\right).
\]
En efecto, la poda deja $\log(1-\theta^2)$ en el ideal del producto de los cuadrados del ciclo.
Los operadores de campo $\zeta_i\partial_{\zeta_i}$ conservan ese ideal.
Para cada $i$ del ciclo se prolongan en la cara tanto $f_{c-1}|_{\zeta_i=0}$ como
$\partial_{\zeta_i}f_{c-1}|_{\zeta_i=0}$, ambas nulas como gérmenes allí.
La serie de Taylor en el polidisco sólo admite así monomios con todas las multiplicidades
indicadas. Una identidad comprobada sólo al evaluar las activaciones físicas no bastaría.
En el bosque los grados menores que $c+1$ se anulan. Con varios ciclos sólo se usa
la anulación por debajo de $c-1$; el coeficiente de ese orden puede sobrevivir.

\paragraph{Cotas de clase y orden de elección.}
Sean $R\asymp1/v$ y $r_i\asymp1/d_i$ radios interiores fijos en proporción a los de arriba.
Póngase $z=|s|/R$ y $t_i=|\zeta_i|/r_i$. Schwarz sucesivo en cada campo da el
factor $\prod_i t_i$ y, para $f_{c-1}$ unicíclico, otro factor $\prod_{i\in\mathcal C}t_i$.
Cauchy en $s$ y la serie geométrica acotan la cola que empieza en $w$ por
$C^v\prod_i t_i\, z^w/(1-z)$. Se usa $w=c+1,c,c-1$ para las tres colas.
El término aislado del ciclo está acotado por
$C^v z^{c-1}\prod_i t_i\prod_{i\in\mathcal C}t_i$.
En el punto físico $z\le Cv/j$, $t_i\le C d_i/j$. Cuando $\max t_i\le1/2$,
$\prod_{i\in\mathcal C}t_i\le\sum_{i\in U}t_i\le Cv^2/j$.
Se obtienen los tres exponentes enunciados, absorbiendo ese $v^2$ en el factor indicado.

Los radios, la cota física de los campos y $C$ son de clase, anteriores a $\rho$ y $b$.
Después de fijar $b$, se elige $n_0(b)$ para que $Cv/j$ y $Cd_i/j$ quepan en los
radios relativos, y $Cq\le\sqrt j$ en todo el intervalo. Ello usa $d_i\le v\le8\log n$.
Para el reemplazo posterior por el valor exacto, el error por cara es a lo sumo
$e^{Cv}((Cq/j)^q+Ce^{-c_0j})$; la alternancia cuesta a lo sumo $2^v$ por soporte.
La cota sumada de átomos físicos es la de T2.2 de R21, con su estimación más fuerte
$Cq\le j^{1/3}$ al sumar los soportes. La continuación en $s$ se ha aplicado únicamente
al aproximante. Todo el argumento se hace en una máscara fija antes de integrar en $\nu$.

\paragraph{Referencias de control.}
Acta de jets, puntos 1a--1g, 2e--2g y 3; especialmente líneas 145--165 y 235--260.
X03 es la parte 03 de la entrega de cotejo de la adenda 2, de 17 de noviembre, 07:21 UTC.
Las rutas de X03, de B03 y de los ``Grados intermedios'' figuran en la lista de precisiones
de la ronda 22, en la sección sobre su ubicación:
\path{velocidad_via2_ronda22_adenda_torres_20261117T133158Z/}\allowbreak{}%
\path{PRECISIONES_PARA_EL_CONSOLIDADO.md}. Allí X03 es
\path{velocidad_via2_cotejo_adenda2_20261117T072100Z/latex/03_truncacion.tex}, líneas 2--34.
Las demás claves de las entradas son las del catálogo literal.
